Everything presented as output in this guide is regenerable from what ships in
the repository, in the QtRocket documentation tradition of docs whose claims
can be re-derived rather than trusted.

\paragraph{Transcripts.} The five tutorial scripts live in
\srcf{docs/qtrocket\_userguide/scripts/} and run verbatim:
\code{qtrocket-cli tutorial1.cli} (from a directory containing the data files
each tutorial's setup box names). The typeset transcripts under
\srcf{docs/qtrocket\_userguide/transcripts/} were generated by running each
script through the release binary and interleaving input with captured output
--- for every command, the script prefix up to that command was run in a fresh
process and its output diffed against the previous prefix's, so each command
is paired with exactly the lines it produced. Tutorial~3 exits with code~1 by
design (it contains two intentional failures).

\paragraph{Screenshots.} The four figures are \code{qtrocket-visualizer}
displaying the \code{.qrd} files the tutorial sessions saved, captured at
$1200\times800$ under a virtual X server (Xvfb, llvmpipe software GL).
Tutorial~3's pair uses the Hi-Vis color scheme so the resolver's fixed error
red is unmistakable against the orange nose and white body.

\paragraph{Plot data.} The \code{data/*.dat} files are two-column extracts of
the \repl{save}d trajectory CSVs --- the $(t,z)$ columns, or $(x,z)$ for the
tilted flight --- downsampled to roughly 150 rows (final row kept) and
whitespace-separated for pgfplots, so the figures rebuild from committed data
without rerunning simulations. The values themselves are unmodified.

\paragraph{This PDF.} \code{make} in \srcf{docs/qtrocket\_userguide/} runs
\code{latexmk -pdf -shell-escape} (minted needs shell escape). The style is
the QtRocket documentation series' shared identity --- XCharter text,
Inconsolata code, the \code{qr*} palette --- inherited from the whitepapers
under \srcf{docs/}.

\paragraph{Versions.} Captured against QtRocket v0.2
(\code{development}, July 2026), motor data from the bundled
\srcf{data/Aerotech.rse} and \srcf{tests/data/motors/estes\_small.qmd}.
Simulations are deterministic; rerunning the scripts reproduces every number
shown, including the CSV contents.
